\documentclass[aspectratio=169,usepdftitle=true,11pt,ngerman,t]{beamer}

\input{../global/user}
\errorcontextlines999999

% base
\usepackage[utf8]{inputenc}
\usepackage[T1]{fontenc}
\usepackage{lmodern}
\usepackage{microtype}
\usepackage[main=ngerman]{babel}
\usepackage{csquotes}

% graphics
\usepackage[prefix=]{xcolor-material}
\usepackage{graphics}
\usepackage{tcolorbox}
\usepackage{twemojis}
\usepackage{tikz}
\usetikzlibrary{
    positioning,
    calc,
    tikzmark
}

% math
\usepackage{amsmath,amssymb}
\usepackage{mathtools}
\usepackage{array}
\usepackage{booktabs}
\usepackage{cancel}
\usepackage{ulem}

% flo
\usepackage{fancyqr}
\usepackage[addons]{color-palettes}
\usepackage[encoding,noloadlangs,cpalette,numinpar,fakeminted]{sopra-listings}
\usepackage{code-animation}
\usetheme[theorems]{colorful-dream}
\theoremstyle{exercise}
\newtheorem*{exercise*}{Aufgabe}
\theoremstyle{solve}
\newtheorem*{solve*}{Lösung}

% tools
\usepackage{algpseudocode}
\usepackage{dirtree}

% warnings
\usepackage{silence}
\WarningFilter{latexfont}{Some font shapes were not available}
\WarningFilter{latexfont}{Font shape}
\let\origenddocument\enddocument\renewcommand{\enddocument}{\hfuzz=10.95pt\origenddocument}
% title image: box
% ([vshift], inverse scale, width, content)
\newcommand\titleimagebox[4][0pt]{%
\newlength\tibheight\setlength\tibheight{4cm}
\newlength\tibtY\setlength\tibtY{-6em}
\titleimage{\begin{minipage}[c][#2\tibheight]{#3}\vspace*{#2\tibtY}%
    \vspace*{#1}\begin{tcolorbox}[boxrule=0.1pt]%
    \centering#4% TODO: bug: not centered if too wide
    \end{tcolorbox}%
\end{minipage}}}

% outro image: text
% (text)
\newcommand\outroimagetext[1]{\outro{Karlsruhe, \today%
\begin{tikzpicture}[remember picture,overlay]
    \node[align=center] at (current page.center) {#1};
\end{tikzpicture}}}% chktex 31

% reset counters
\newcommand\resetframecounters{\addtocounter{exercise}{-1}\addtocounter{solve}{-1}}
\newcommand\resetsolve{\addtocounter{solve}{-1}}

% small footnote
\newcommand{\nicefootnote}[1]{\tikz[overlay,remember picture]{\node[anchor=south,text width=\dimexpr\textwidth+2em\relax,align=center,above=1em of current page.south] (nicefootnote) {\hrule\vspace{1pt}\footnotesize#1};}}
\title{Digitaltechnik}
\institute{Karlsruher Institut für Technologie}

\author{\lqfullname}
\email{\lqemail @student.kit.edu}% chktex 1
\input{../global/credit}

%%%%%%%%%%%%%%%%%%%%%%%%%%%%%%%%%%%%%%%%%%%%%%%%%%%%%%%%%%%%%%%%

\subtitle{Tutorium 0}
\date{05. Mai 2026}

\titleimagebox[-100cm]{1}{5cm}{}

\outroimagetext{\LARGE Bis nächste Woche!}

%%%%%%%%%%%%%%%%%%%%%%%%%%%%%%%%%%%%%%%%%%%%%%%%%%%%%%%%%%%%%%%%

\title{\texorpdfstring{\rlap{Digitaltechnik und Entwurfsverfahren}}{Digitaltechnik und Entwurfsverfahren}}

\begin{document}

\title{Digitaltechnik}
\begin{frame}[plain]\titlepage\end{frame}

\section{Willkommen}

\subsection{Tutorium}

\begin{frame}{Seid ihr hier richtig?}\onslide<+->
    \begin{itemize}[<+->]
        \item[\alt<.>{\usebeamertemplate{itemize item}}{\faCheck}] Tutorium am Dienstag um 15:45 Uhr?\pause%
        \begin{itemize}
            \item Beginn: \emph{dritte} Vorlesungswoche
        \end{itemize}
        \item[\faCheck] Für die Vorlesung Digitaltechnik und Entwurfsverfahren?
        \item[\faCheck] Euer Tutor heißt \lqfullname? \onslide<+->\textit{\color{lightgray} (Das bin ich!)}
    \end{itemize}
    \vspace{1em}
    \begin{center}
        \onslide<+->\fontsize{50}{70}\selectfont\texttwemoji{tada}
    \end{center}
\end{frame}

\subsection{Vorstellungsrunde}

\begin{frame}{About me\dots}\onslide<+->
    \begin{itemize}
        \item \lqfirstname%
        \item Master Informatik
        \item mag \LaTeX~und Minecraft~:D
        \vspace{2em}\item<+-> \dots und ihr?
    \end{itemize}
\end{frame}

\section*{Organisation}

\subsection{Übungsblätter}

\begin{frame}{Übungsblätter}\onslide<+->
    \begin{itemize}[<+->]
        \item Einzelabgabe
        \begin{itemize}
            \item Benennung
            \item \emph{Rechtzeitig!}
        \end{itemize}
        \item \emph{handschriftlich}
        \begin{itemize}
            \item Tablet erlaubt
        \end{itemize}
        \item 2 Bonuspunkte
        \begin{itemize}
            \item insgesamt 50\% der Punkte
            \item maximal 2 Blätter nicht abgegeben
            \item \enquote{aktive Mitarbeit} im Tutorium
        \end{itemize}
    \end{itemize}
    \vspace{1em}\onslide<+->
    \begin{quotation}
        Eine Person, die eine korrekt gelöste (Teil-) Aufgabe abgegeben hat, muss auch in der Lage sein, diese im Tutorium vorzurechnen.
    \end{quotation}
\end{frame}

\subsection{Tutorium}

\begin{frame}{Tutorium}\onslide<+->
    \begin{itemize}[<+->]
        \item keine Anwesenheitspflicht\onslide<+->, aber \emph{sinnvoll!}
        \begin{itemize}
            \item deutlich höhere Wahrscheinlichkeit, die Klausur zu bestehen!
        \end{itemize}
        \item Übungsblätter besprechen
        \item Fragen stellen
        \item Aufgaben üben
    \end{itemize}
    \onslide<+->\vspace{1em}
    \begin{block}{Folien auf GitHub}
        \centering\fancyqr[size=2.5cm]{https://github.com/Laqco/kit-dt-ss26}
    \end{block}
\end{frame}

\subsection{Klausur}

\begin{frame}{Klausur}\onslide<+->
    \begin{itemize}[<+->]
        \item normal: März 2027
        \item \textit{\color{lightgray} Einzel- und Nachklausur 14.08.2026, 08:00}
    \end{itemize}
\end{frame}

\section[Das nullte Tutorium]{Tutorium}

\subsection{Fragen und Übung}

\begin{frame}{Hier könnte Ihre Werbung stehen!}\onslide<+->
    \begin{itemize}[<+->]
        \item Fragen zum Übungsblatt?
        \item Übungsaufgaben?
    \end{itemize}
\end{frame}

\end{document}